\section*{Oppgave 4 — Bærekraft}
\setcounter{section}{4}
\addcontentsline{toc}{section}{Oppgave 4 — Bærekraft}

\setcounter{subsection}{0}

\subsection{Kviles bærekraft}

\textbf{Økonomisk} har Kvile et godt utgangspunkt med ca.\ 200~MNOK i omsetning, vekst og en stabil offentlig kunde. \textbf{Sosialt} bidrar bedriften med industriarbeidsplasser i Orkanger og møbler til offentlige rasteplasser. \textbf{Miljømessig} er vurderingen svakere: aluminium er energikrevende, lakkering kan gi kjemikalie- og VOC-utslipp, transporten skjer med lastebil over hele Norge, og RAST-produktene bruker engangsemballasje av glassfiberstrie.

\subsection{Endringsforslag}

Det mest åpenbare tiltaket gjelder emballasjen. Glassfiberstrien rundt RAST-produktene er en ren engangsløsning som kastes ved montering. Kvile bør heller utvikle returemballasje — f.eks.\ vevde tekstilposer eller stålrammer — som sjåførene tar med tilbake etter levering. Siden Kviles egen transportavdeling utfører leveransene, er returlogistikken allerede på plass.

Kvile bør også gå over til resirkulert aluminium, som krever ca.\ 95\,\% mindre energi enn primæraluminium \parencite{iai2025}.

Når det gjelder transport og produktdesign ser vi to muligheter. Fabrikken ligger nær kaia på Orkanger, men Kvile bruker utelukkende lastebil. For leveranser til kystbyer kan sjøtransport kombinert med lokal distribusjon redusere utslipp, og ruteoptimalisering av lastebilturene kan redusere tomkjøring. Samtidig bør Kvile designe produktene slik at slitedeler som treplatene kan skiftes uten å bytte hele møbelet. Dette er spesielt relevant for Vegvesenet, som i dag bestiller komplette erstatninger for slitte møbler.

Flere av disse tiltakene henger sammen med Lean-prinsippene fra oppgave~2. ATO reduserer overproduksjon og unødvendig lager, returemballasje kutter materialavfall, og bedre ruter betyr mindre transportsløsing.